\documentclass[11pt]{article}
\usepackage[margin=1in]{geometry}
\usepackage{amsmath,amssymb}
\usepackage{booktabs}
\usepackage{graphicx}
\usepackage{xcolor}
\usepackage[colorlinks=true,linkcolor=blue!60!black,citecolor=blue!60!black,urlcolor=blue!60!black]{hyperref}
\usepackage{microtype}

\newcommand{\system}{\textsc{LocalFlash}}
\newcommand{\dflash}{DFlash~2}
\newcommand{\omlx}{oMLX}

\title{\textbf{Losing the Draft, Keeping the Speed:\\ Block-Diffusion Speculative Decoding and Prefix-Cached Serving\\ for Long-Context Agentic LLM Workloads on Apple Silicon}}

\author{chengyixu\thanks{Personal-systems research. Deployment artifacts and raw measurements:
\url{https://github.com/chengyixu} (code) and the accompanying Hugging Face repository.
Hardware: MacBook Pro, Apple M4 Max (16-core), 64\,GB unified memory, macOS 27.0.}}
\date{August 24, 2026}

\begin{document}
\maketitle

\begin{abstract}
We document the end-to-end engineering of a single-model local deployment for
coding, engineering, and agentic workloads on an Apple M4 Max workstation with
64\,GB of unified memory, targeting two constraints: native context lengths of
at least 50{,}000 tokens at no worse than 4-bit weight precision. Starting from
a llama.cpp-based deployment of a 27B dense model accelerated by the \dflash{}
block-diffusion drafter---which achieved only 9--15 tokens/s decode and required
435 seconds of time-to-first-token (TTFT) per 32k-token turn---we evaluate four
serving stacks under identical workloads: two llama.cpp configurations, vanilla
MLX inference, and a community fork of the \omlx{} server integrating \dflash{}
through a dedicated block-diffusion engine. The winning configuration couples
the official 4-bit MLX quantization of Qwen3.8-27B with the \dflash{} drafter
(quantized to 4 bits at load time, block size 5, adaptive verification) and a
paged prefix cache. It sustains a median decode rate of 809 tokens/s on
short-context reasoning (a $54\times$ improvement over the baseline), reduces
per-turn TTFT in multi-turn agentic sessions from 435\,s to 8--16\,s after the
initial prefill ($26\times$), preserves lossless output, and passes long-context
retrieval probes at 70k tokens (5/5 ordered needle recall) alongside deterministic
code-generation sanity checks. We report the full methodology, negative results,
measurement pitfalls (including a silent memory-contention failure mode that
invalidated an entire benchmark matrix), and release all harness code, raw
measurements, and deployment manifests for reproducibility.
\end{abstract}

\section{Introduction}
Local inference has reached an inflection point: open-weight dense models in the
27B class now fit comfortably within the unified memory of high-end consumer
hardware at 4-bit precision, while remaining competitive with frontier systems
on software-engineering benchmarks. Qwen3.8-27B~\cite{qwen38} reports
73.0 on Terminal-Bench 2.1 (Terminus) and 61.7 on SWE-bench Pro, and its hybrid
architecture---64 layers alternating gated DeltaNet linear attention with gated
full attention (GQA, 24 query / 4 key-value heads)---yields an unusually small
KV footprint, making very long contexts practical on unified-memory machines.

This paper is a systems case study with three contributions:

\begin{enumerate}
\item \textbf{A measured comparison of four serving stacks} for the same target
model on Apple Silicon, spanning GGUF/llama.cpp and MLX ecosystems, with and
without speculative decoding (Section~\ref{sec:results}).
\item \textbf{An agentic-serving evaluation protocol} that measures what coding
agents actually experience: cold-context prefill once, followed by many
incremental turns whose prompts extend a shared, growing context. This protocol
exposes a $26\times$ TTFT gap that single-turn benchmarks hide
(Section~\ref{sec:agentic}).
\item \textbf{A documented measurement methodology}, including calibration of
synthetic prompts against the tokenizer, integrity verification of downloaded
weights, and a post-mortem of an invalidated benchmark run caused by
memory contention among concurrent servers (Section~\ref{sec:method},
Section~\ref{sec:pitfalls}).
\end{enumerate}

All artifacts---benchmark harnesses, deployment controller, per-model engine
configuration, raw JSON-lines measurements---are released alongside this paper.\footnote{Code:
\url{https://github.com/chengyixu/qwen38-dflash2-bench}; model configuration card:
Hugging Face, same identifier.}

\section{Background}
\label{sec:background}

\subsection{Speculative decoding and block diffusion}
Speculative decoding accelerates autoregressive generation by drafting multiple
tokens with a cheap model and verifying them in parallel against the target;
acceptance is defined so the output distribution is preserved exactly
(lossless decoding)~\cite{leviathan2023fast}. Classical drafters propose tokens
sequentially, capping speedups near $2$--$3\times$. \dflash{}~\cite{dflash}
instead drafts with a lightweight \emph{block-diffusion} model that predicts an
entire block in one pass and keeps top candidates at every position; a selector
traces one coherent path, and two-tap dynamic convolutions prevent quality decay
toward the end of the block. On an H200, \dflash{} reports acceptance lengths of
up to 5.46 tokens per verification step and $3.43\times$ throughput over
autoregressive decoding on GSM8K~\cite{dflash2card}.

\subsection{Agentic serving and prefix reuse}
Coding-agent workloads differ fundamentally from chat: each tool-call turn
re-submits a context that is a strict extension of (or a bounded edit to) the
previous one. FreeToken~\cite{freetoken} demonstrates that serving systems
which preserve recurrent state across context edits reduce worst-case TTFT from
minutes to seconds, and that tail latency above ${\sim}150$\,s crosses real
client timeout thresholds. On Apple Silicon, \omlx{}~\cite{omlx} implements the
analogous idea as a paged, hot/cold KV-cache with prefix sharing, reporting
agent TTFT reductions from 30--90\,s to under 5\,s for repeated prefixes.

\subsection{Target model}
Qwen3.8-27B~\cite{qwen38} is a vision-capable dense model: hidden size 5120,
64 layers in a $16\times(3\times(\text{DeltaNet}\!\to\!\text{FFN}) \to
1\times(\text{Gated Attention}\!\to\!\text{FFN}))$ layout, 248{,}320-vocabulary,
native 262{,}144-token context (YaRN-extensible to 1M), trained with multi-token
prediction heads and shipped with official \texttt{reasoning\_effort} control.
Its recommended thinking-mode sampling is temperature~$1.0$, top-$p$~$0.95$,
top-$k$~$20$.

\section{Experimental Setup}
\label{sec:method}

\subsection{Hardware and software}
All experiments ran on one machine (Table~\ref{tab:env}). The GPU wired-limit
advisory reported by \omlx{} was left at Apple's default (51.8\,GB working-set
cap); the process-level memory ceiling observed during serving was 42.3\,GB.

\begin{table}[h]
\centering
\caption{Environment.}
\label{tab:env}
\small
\begin{tabular}{ll}
\toprule
Component & Version \\
\midrule
Machine & MacBook Pro, Apple M4 Max (16 CPU cores), 64\,GB LPDDR \\
OS & macOS 27.0 (26A5406e) \\
MLX / mlx-lm & 0.32.1 / 0.31.3 (Python 3.14 venv) \\
\omlx{} z-lab fork & 0.6.2-dflash2 (\texttt{dflash\_mlx} 0.1.10+omlx.5) \\
llama.cpp (baseline) & bundled DFlash2 build (retired profile) \\
llama.cpp (upgraded) & 0.2.0, build 10566 \\
Weights & \texttt{mlx-community/Qwen3.8-27B-4bit} @ \texttt{3e6447f} \\
Drafter & \texttt{z-lab/Qwen3.8-27B-DFlash2} (BF16, 2B params) \\
\bottomrule
\end{tabular}
\end{table}

Weight integrity was enforced by SHA-256 comparison against the Hugging Face
LFS manifest. This caught two corrupted shards (Section~\ref{sec:pitfalls});
all final artifacts verify cleanly. Model discovery on the serving path exposes
three entries (target, drafter, and a separate MTP-head adapter); only the
target is user-facing, with the drafter attached transparently by the engine.

\subsection{Engine configuration}
The winning profile enables \dflash{} for the target model with: drafter
quantized at load time to 4-bit weights (group size 64, 16-bit activations),
block size 5 (following z-lab's guidance for quantized targets, since MLX's
quantized matmul loses efficiency at wider verify shapes), adaptive
verification, sliding-window drafting at 2048 tokens, and an in-memory L1
prefix cache capped at 8\,GiB across 4 entries. SSD-tier caching was left
disabled for these experiments.

\subsection{Workloads and metrics}
Synthetic workloads are generated from a fixed technical vocabulary and
calibrated against the target tokenizer: the word-to-token ratio converges to
$1.483 \pm 0.001$ across 4k--74k targets, so prompt sizes are specified in
true tokens. Each request streams a chat completion; we record
\begin{itemize}
\item \textbf{TTFT}: request start to first streamed content delta;
\item \textbf{Prefill throughput}: provider-reported prompt tokens $\div$ TTFT;
\item \textbf{Decode throughput}: completion tokens beyond the first, divided
by the post-TTFT interval;
\item \textbf{Quality gates}: (i) arithmetic probes embedded after long filler
context (expected answer computable without reading the filler);
(ii) five ``needle'' records planted at even depths, requiring exact ordered
recall; (iii) a code-generation probe requiring a specific function signature,
guard clauses, and an $O(n)$ structure.
\end{itemize}
Sampling follows the model card for thinking mode ($T{=}1.0$, top-$p$~$0.95$)
for performance runs; quality probes use greedy decoding. Unless noted,
short-context results are medians over 3 runs after a warm-up request.

\subsection{Tested stacks}
\label{sec:stacks}
\begin{enumerate}
\item[\textbf{A.}] \textbf{Baseline}: retired llama.cpp build with
GGUF Q4\_K\_M finetune + \dflash{} drafter GGUF, 32k context, served via
LaunchAgent (the incumbent deployment).
\item[\textbf{B.}] \textbf{Upgraded llama.cpp}: build 10566, same GGUF pair;
the drafter failed to load (tensor-layout mismatch), so stack B degenerates to
autoregressive decoding at 64k context.
\item[\textbf{C.}] \textbf{Vanilla MLX}: \texttt{mlx-lm} direct generation on
the 4-bit MLX checkpoint (no server features).
\item[\textbf{D.}] \textbf{\omlx{} z-lab fork + \dflash{}} (\system{}):
the proposed stack (Section~\ref{sec:stacks-config}).
\end{enumerate}

A fifth configuration---mainline \omlx{} with an MTP head---was evaluated early
and rejected: singleton MTP verification through the compatibility patch was
far slower than \dflash{}, and mainline lacks the block-diffusion engine.

\section{Results}
\label{sec:results}

\subsection{Single-turn benchmarks}
Table~\ref{tab:main} summarizes single-turn performance. Two results dominate.
First, \emph{decode}: \dflash{} under the fork reaches a median 809 tokens/s on
short contexts versus 15 tokens/s for the incumbent deployment ($54\times$),
consistent with acceptance-length dynamics rather than hardware change.
Second, \emph{long-context prefill}: no stack exceeds ${\sim}145$ prompt
tokens/s at 24k--33k tokens; fresh-prompt TTFT remains minutes-scale for every
stack. Prefill, not decode, is the binding constraint for local long-context
use---motivating Section~\ref{sec:agentic}.

\begin{table}[h]
\centering
\caption{Single-turn results on identical synthetic workloads. Decode values
are medians over 3 short-context runs (265-token prompt, 384 requested tokens;
generations terminate early under thinking mode). $^{*}$Stack~B's drafter could
not load; its decode is plain autoregressive. $^{\dagger}$Vanilla MLX figures
from direct generation (68-token probe).}
\label{tab:main}
\small
\begin{tabular}{lcccc}
\toprule
Stack & Decode (tok/s) & TTFT@32.5k (s) & Prefill@32.5k (tok/s) & Context cap \\
\midrule
A. llama.cpp+\dflash{} (old) & ${\sim}15$ (smoke) & 435.5 & 74.5 & 32k \\
B. llama.cpp AR (new)$^{*}$ & 352.9 & 373.0 & 87.0 & 64k \\
C. Vanilla mlx-lm$^{\dagger}$ & ${\sim}10$ & --- & --- & --- \\
D. Fork+\dflash{} (\system{}) & \textbf{809.1} & 486.0 & 66.8 & 262k \\
\bottomrule
\end{tabular}
\end{table}

At intermediate context (5{,}487 tokens), stack D records TTFT 38.1\,s
(144.1 prompt tokens/s) and 613.6 tokens/s decode; the non-monotonic prefill
rate across context lengths reflects chunked-prefill scheduling in the fork's
engine rather than instability.

\subsection{Multi-turn agentic protocol}
\label{sec:agentic}
We simulate a coding-agent session: turn~1 submits a 32.5k-token document plus a
question; turns~2--3 append short follow-ups whose prompts extend the shared
prefix (32{,}513 and 32{,}549 tokens respectively). Under stack A, \emph{every}
turn pays the full re-prefill (${\sim}435$\,s). Under stack D the L1 prefix
cache absorbs turns 2--3:

\begin{table}[h]
\centering
\caption{Agentic session on stack D (32.5k base context). Turn-1 cost is paid
once; subsequent turns are $26$--$52\times$ faster than the incumbent's
per-turn cost.}
\label{tab:agentic}
\small
\begin{tabular}{llrrl}
\toprule
Turn & Kind & Prompt (tok) & TTFT (s) & Note \\
\midrule
1 & cold base & 32{,}478 & 428.08 & full prefill \\
2 & cached prefix & 32{,}513 & \textbf{8.24} & $+35$ tokens \\
3 & cached prefix & 32{,}549 & \textbf{15.78} & longer generation \\
\bottomrule
\end{tabular}
\end{table}

This reproduces on unified memory the qualitative finding of
FreeToken~\cite{freetoken} on discrete GPUs: prefix/state reuse, not kernel
speed, governs perceived agentic latency. Tail risk also collapses: the
incumbent's every-turn 435\,s sits far above client watchdog thresholds,
whereas stack D's steady-state turns stay below 16\,s.

\subsection{Quality gates}
Long-context competence was probed at a true 70{,}249 tokens (five needles at
depths $1/6$--$5/6$): the model recalled all five records \emph{in order}
(5/5) with greedy decoding. An earlier 33.8k-token probe scored 4/5 solely due
to a 100-token response cap truncating the list---a harness artifact we fixed
by raising the cap to 600, underscoring that reasoning-mode outputs must budget
for self-generated chain-of-thought before the visible answer. Code-generation
probes passed deterministically (correct signature, guard clauses returning
empty lists for invalid windows, accumulator-based $O(n)$ body). Because
\dflash{} verification is lossless, these quality results transfer unchanged
from the underlying target model; we did not observe any speculative-path
artifacts across hundreds of generations.

\section{Measurement Pitfalls and Negative Results}
\label{sec:pitfalls}

\paragraph{Silent memory contention invalidates matrices.} Our first two
``clean-environment'' benchmark batches against the fork were unknowingly
directed at a stale mainline \omlx{} instance that had auto-respawned via a
package-manager service and bound the target port first; meanwhile a llama.cpp
instance held ${\sim}18$\,GB and system swap grew past 26\,GB. Results looked
plausible yet were $3$--$30\times$ off (decode collapsing to 19 tokens/s at
24k context). Detection required cross-checking process identity, port
ownership, and swap pressure. We now require, before every batch: exactly one
listening process on the port, positive confirmation of the engine banner in
server logs, zero competing inference processes, and stable swap. We publish
the invalidated numbers' lesson but exclude the runs from all tables.

\paragraph{Downloaded weights can corrupt plausibly.} Resuming interrupted
transfers across two different downloaders produced shards with valid
safetensors headers and correct file sizes but corrupted payloads; the model
loaded ``successfully'' and emitted fluent multilingual gibberish at temperature
zero. Header checks passed; only SHA-256 against the hub manifest revealed the
damage. Any reproducibility effort must hash-verify weights, not merely load
them.

\paragraph{Draft-format coupling is fragile.} Stack B's inability to load the
\dflash{} GGUF (``wrong number of tensors; expected 81, got 58'') shows that
speculative pairs must be built for the consuming runtime's tensor layout; a
drafter that works in one engine is not portable to another major version.

\paragraph{Quantized verify widths matter.} Following z-lab's guidance we set
block size 5 for the 4-bit target; larger verify blocks degrade MLX's quantized
matmul efficiency, which would silently erode the acceptance-length gains.

\paragraph{Response budgets interact with thinking mode.} Quality probes must
budget for the model's own reasoning tokens; caps sized for answers alone
manufacture false failures (Section~\ref{sec:results}).

\section{Deployment}
\label{sec:deployment}
The production profile serves \texttt{Qwen3.8-27B-4bit} on
\texttt{127.0.0.1:7870} behind a KeepAlive LaunchAgent with a controller
supporting \texttt{status/start/stop/restart/test/logs}; status distinguishes
\textsc{starting} (model loading into Metal, ${\sim}60$\,s cold) from
\textsc{running} (health-checked). The endpoint exposes OpenAI-compatible
chat/completions plus an Anthropic-compatible messages surface, with bearer
authentication. Client integrations were migrated to report the true context
window (262k) and API type. Disk footprint: 16.6\,GB (target) + 3.4\,GB
(drafter); the superseded 18\,GB GGUF pair was retired after validation.

\section{Related Work}
Beyond the systems above~\cite{dflash,dflash2card,freetoken,omlx,qwen38}, our
measurements echo classical speculative-decoding theory~\cite{leviathan2023fast};
paged attention~\cite{pagedattention} underlies the fork's cache design; and the
gated DeltaNet layers~\cite{gdn} in Qwen3.8 explain its small full-attention KV
share. Community ports (dflash-mlx inside the fork; MTPLX-style verify kernels)
were instrumental; none of the CUDA-native edge-MoE systems~\cite{freetoken}
run on Apple Silicon, which motivates documenting the MLX-side path.

\section{Limitations and Future Work}
Measurements reflect one machine and synthetic contexts; real repository
contexts may stress the drafter differently (acceptance length typically falls
on code-heavy distributions~\cite{dflash2card}). Single-run long-context cells
carry thermal-state variance; the agentic deltas ($26\times$) are far outside
that noise band, but absolute TTFT values should be treated as indicative.
Future work: SSD-tier prefix eviction studies across days-long agent sessions;
acceptance-length instrumentation per workload class; 6-bit target comparisons
within the 42\,GB ceiling; and controlled concurrency scaling beyond the
single-stream focus here.

\section{Conclusion}
A careful re-assembly of existing pieces---a 4-bit MLX quantization, a
block-diffusion drafter, and a community serving fork with prefix caching---
converts a barely usable local deployment (15 tokens/s, minutes-per-turn,
32k context) into a genuinely agentic one (809 tokens/s median decode,
seconds-per-turn after one prefill, 262k native context, lossless output) on a
single M4 Max. We release everything needed to reproduce or adapt it.

\begin{thebibliography}{9}
\bibitem{qwen38} Qwen Team. \emph{Qwen3.8-27B Model Card}. Hugging Face, 2026.
\bibitem{dflash} J.~Chen, Y.~Liang, Z.~Liu. \emph{DFlash: Block Diffusion for Flash Speculative Decoding}. arXiv:2602.06036, 2026.
\bibitem{dflash2card} Inco AI. \emph{DFlash 2: Keep Drafting Parallel}; \texttt{z-lab/Qwen3.8-27B-DFlash2} model card, 2026.
\bibitem{freetoken} S.~Yang et al. \emph{FreeToken: Efficient Edge-Native MoE Serving with Bandwidth-Adaptive Execution}. arXiv:2608.16157, 2026.
\bibitem{omlx} J.~Kim. \emph{oMLX: LLM Inference Server with Continuous Batching \& SSD Caching for Apple Silicon}. GitHub, 2026.
\bibitem{leviathan2023fast} Y.~Leviathan, M.~Kalman, Y.~Matias. \emph{Fast Inference from Transformers via Speculative Decoding}. ICML, 2023.
\bibitem{pagedattention} W.~Kwon et al. \emph{Efficient Memory Management for Large Language Model Serving with PagedAttention}. SOSP, 2023.
\bibitem{gdn} S.~Yang, J.~Kautz, A.~Hatamizadeh. \emph{Gated Delta Networks: Improving Mamba2 with Delta Rule}. arXiv:2412.06464, 2024.
\end{thebibliography}

\appendix
\section{Artifact Inventory}
\texttt{bench/bench\_api.py} (calibrated workload generator + streaming
benchmarker), \texttt{bench/qa\_check.py} (needle/code gates),
\texttt{bench/agent\_turns.py} (prefix-reuse protocol),
\texttt{bench/hf\_fetch.py} (resumable, retry-hardened downloader),
\texttt{deploy/local-llm.command} (controller),
\texttt{deploy/omlx-model-settings.example.json} (engine config),
\texttt{results/raw/*.jsonl} (39 records: 32 benchmark rows, 3 agentic turns,
4 quality-gate outcomes). Secrets are redacted; the API key appears nowhere in
released files.

\section{Reproduction Steps}
Install \omlx{} fork 0.6.2-dflash2; download the two repositories listed in
Table~\ref{tab:env}; place the example settings at
\verb|~/.omlx/model_settings.json|; start the server; run
\verb|python bench/bench_api.py --base-url http://127.0.0.1:7870 ...| with the
labels used in Table~\ref{tab:main}; verify single-instance ownership of the
port before each batch (see Section~\ref{sec:pitfalls}).

\end{document}
